%!TEX root = ../../csuthesis_main.tex
\chapter{绪论}
\section{研究背景与意义}
\subsection{研究背景}
随着具身智能、数字人仿真、生物力学仿真技术快速发展，基于物理引擎的仿人运动生成已经成为机器人、虚拟现实、康复工程领域的研究热点。传统仿人控制方案大多采用力矩驱动人形机器人模型，仅关注关节层面运动输出，忽略人体真实生理特征，难以复现人体肌肉‑肌腱的非线性驱动、关节生物约束、肌肉冗余驱动等生理特性，生成运动常常出现动作失真、不符合人体运动学规律等问题arXiv。

运动捕捉（MoCap）数据是人体运动生成的重要数据源，但直接复刻运动捕捉数据的方法泛化能力差，无法响应高级指令，很难完成未见过的目标任务。强化学习为物理仿真环境下的运动模仿提供可行路径，现有很多强化学习人形控制框架侧重于关节力矩控制，缺少对肌肉骨骼系统完整建模，很难同时输出合理关节姿态与符合生理的肌肉激活信号，和真实人体肌电 EMG 测量结果匹配度较低，制约了数字人、康复仿真、运动生物力学的进一步落地应用。

KINESIS 开源框架的出现为肌肉骨骼人体运动模仿提供新的基础工具。该框架基于无模型强化学习，使用约 1.8 小时人体运动捕捉 locomotion 数据集训练，借助负样本挖掘得到人体运动先验，能够驱动最高 290 块肌肉的复杂肌肉骨骼模型；同时支持下游任务：文本到运动控制、目标点导航、专项运动动作生成，输出肌肉激活模式和真实人类 EMG 肌电信号具备相关性，为生理可信的人体运动仿真提供了开源基线系统arXiv。

\subsubsection{研究意义}
理论意义：
第一，深化肌肉骨骼系统强化学习运动模仿的相关研究。传统人形仿真多以力矩控制为核心，本研究基于 KINESIS 肌肉驱动范式，探究**肌肉冗余驱动条件下运动模仿策略学习机制**，对理解 Bernstein 运动冗余问题提供仿真实验参考，丰富生物启发式具身智能的研究案例arXiv。
第二，拓展开源 KINESIS 框架的研究边界。原始论文重点验证行走类运动模仿，本课题在此基础上探究模型在拓展动作、条件生成场景下的性能，补充该开源框架的实验数据与评估结论，完善肌肉骨骼运动模仿领域的开源基线研究资料。
第三，建立兼顾运动姿态精度与生理肌肉激活合理性的综合评估思路，不只看关节角度误差，同时引入肌电相关性指标，为生理可信人体运动生成研究提供评估参考思路。

实际意义：
第一，数字人与虚拟现实方向：优化后的框架可生成生理可信人体全身运动，用于虚拟数字人动画，减少人工动画制作成本，动作更加贴合真实人体运动规律。
第二，康复与运动生物力学仿真：能够输出仿真肌肉激活信号，可以辅助康复动作分析、运动损伤机理仿真，为康复训练模拟提供仿真工具。
第三，人形机器人参考：肌肉驱动仿真得到的运动策略，可以为人形机器人仿人运动控制提供生物启发参考。

\section{国内外研究现状}

\subsection{国内研究现状}
国内在仿人机器人、数字人生成领域成果丰富：一方面，大量研究聚焦于视觉驱动的运动捕捉、基于生成对抗网络、扩散模型的人体姿态序列生成，大多属于运动学层面生成，缺少物理动力学、肌肉骨骼生理约束。

在生物力学仿真方面，国内学者大多基于 OpenSim 开展康复、运动学分析，但是将 OpenSim 和深度强化学习结合实现端到端运动模仿的工作相对偏少。国内对于**肌肉驱动的物理仿真人形强化学习**研究整体起步晚，公开可用的开源基线工具较少，很多研究直接复用国外开源仿真环境。

针对 KINESIS 框架，国内目前直接基于该开源项目展开二次优化、拓展任务的公开研究成果较少，大部分研究停留在力矩人形控制，较少关注仿真肌肉激活与真实肌电信号对齐这一维度。部分高校团队开始关注具身智能的生物真实性，但是缺少可直接复用的优化工程实现。
\subsection{国外研究现状}
在物理仿真人形运动模仿领域，国外起步较早。早期研究主要使用运动捕捉数据模仿，通过模仿学习、强化学习在物理仿真环境复现人体动作。OpenAI、DeepMind 等机构推出多个力矩驱动人形强化学习控制框架，可以实现跑步、行走等多种动作，但模型为力矩驱动，不考虑真实人体肌肉解剖结构，无法输出肌肉激活信息。

在**肌肉骨骼仿真 + 强化学习**方向，国外已经开发 OpenSim 等经典生物力学仿真平台，可以高精度构建人体肌肉骨骼模型，但传统 OpenSim 求解速度慢，较少直接结合端到端强化学习做运动模仿。近年出现一批结合强化学习与肌肉骨骼模型的工作：MyoSuite 系列仿真环境搭建肌肉驱动人形仿真，支持强化学习训练；KINESIS（2025）作为最新开源工作，解决大规模肌肉骨骼模型下无模型运动模仿训练难的问题，使用负挖掘策略提升策略鲁棒性，实现文本条件零样本运动控制，并且验证仿真肌肉输出和真实 EMG 肌电相关性，是该方向重要开源成果arXiv。

同时文本‑to‑motion 生成方向，国外大量生成模型可以根据文本生成人体运动序列，但大多属于 kinematics 纯运动学生成，不经过物理仿真，不考虑肌肉与动力学约束；KINESIS 将文本运动生成和物理肌肉骨骼仿真结合，属于比较新颖的研究路线。
\subsection{研究述评}
综合国内外研究可以看出，肌肉骨骼运动模仿已经完成了从“能否控制肌肉驱动模型”到“能否生成生理可信动作”的跨越，但面向全身协同与实际复现仍有四个方面的不足，构成本课题的切入点。
第一，上下肢协同存在结构性缺口。主流肌肉驱动运动模仿的工作以地面 locomotion 为核心，上肢或未被建模，或仅随躯干被动运动，KINESIS 的作者亦将上肢控制的缺失列为限制动作库多样性的主要因素[1]。缺失协同约束的直接后果是：当任务需要手与脚在时间与空间上配合时，末端执行部位的轨迹只能由躯干摆动间接带动，容易出现末端漂移、接触时刻错位与躯干—肢体时序不一致，动作虽在关节层面“跟得上”，却在任务层面“做不到”。
第二，控制精度与输出稳定性不足。过驱动与高维动作空间使策略输出对奖励权重与观测构成较为敏感，表现为关节轨迹的高频抖动与稳态跟踪误差偏大；在分布外的文本驱动动作上，关键点误差与成功率相对训练分布明显下降。现有工作多通过增加训练数据或加大网络容量应对，较少从观测空间设计、奖励项平滑与采样策略三个层面协同优化精度。
第三，工程复现门槛较高。框架依赖仿真环境、肌骨模型资产与动作数据集的多层耦合，训练流程涉及负样本挖掘的多次迭代与专家权重管理，日志、可视化与评测脚本分散在不同位置；单个专家的训练耗时长，缺少面向有限算力的训练配置建议。这使复现者往往止步于环境配置阶段，难以在基线上开展改进实验。
第四，评价维度较为单一。既有评价以关节角误差、关键点误差与任务成功率为主，肌肉层面则以与肌电信号的整体相关性衡量，缺少对协同一致性、轨迹平滑度与肌肉激活稳定性的量化指标。评价口径不统一，也使得不同改进方案之间难以横向比较。

\section{研究内容与创新点}
本课题以开源项目 KINESIS 为基础，面向生理可信的全身运动生成，完成源码研读、环境复
现、方法改进与对比实验验证，构建具备手脚协同能力、控制精度更高、工程可复现的肌肉骨骼运
动模仿系统。主要研究内容分为四个部分。
\begin{enumerate}
	\item 框架原理梳理与基线复现。研读原始文献与工程代码，拆解整体链路：动作捕捉数据的	筛选与预处理、负样本挖掘的迭代训练机制、策略网络与混合专家结构、肌骨模型与仿真接口、文本条件与高层任务的控制方式；解决依赖安装、模型资产与训练配置等复现问题，完成下肢locomotion 基线复现，并记录原生方案在全身协同动作上的表现缺陷、跟踪误差分布与训练开销。
	\item 面向手脚协同的控制策略改进。原生框架以上肢不受控为前提，缺少肢体间的耦合约
	束。本研究在现有肌骨模型上启用上肢驱动，并在奖励体系中新增协同约束项：以上肢末端的位置与速度跟踪约束其轨迹，以手脚接触与摆动的相对时序约束其节奏，以躯干姿态约束维持整体平衡，从而构建手脚协同的肌肉骨骼运动控制方案，使行走中伸手、跨步搬运、站立取物等复合动作能够被完整生成。
	\item 控制算法精度优化。针对关节轨迹抖动、稳态跟踪误差偏大与肌肉激活输出波动的问题，从三个层面协同改进：优化观测空间的构成与归一化方式，减少冗余与噪声输入；重构多目标加权奖励函数，引入姿态误差的平滑惩罚与权重自适应调节，抑制策略输出的高频振荡；改进负样本挖掘的采样与筛选逻辑，优先保留边界样本与难样本，提升有限训练预算下的收敛质量。
	\item 工程适配、对比实验与评测。搭建统一评测流程，以原生 KINESIS 为基线设置对照实验与消融实验，从关节轨迹误差、动作完成成功率、手脚协同一致性、轨迹平滑度以及肌肉激活与肌电相关性等维度比较改进前后的效果；完成框架的模块化改造、依赖适配与日志可视化封装，整理可复用的工程版本并分析方案的适用边界。
\end{enumerate}
本课题的创新定位于对已有开源框架的功能补全、性能优化与工程落地，不涉及全新框架的提
出，具体包括以下三点。
\begin{enumerate}
	\item 手脚协同控制。原生 KINESIS 以地面运动为核心，上肢不参与主动控制，因而无法表达需要四肢配合的复合动作。本课题在肌肉驱动运动模仿框架内引入上肢末端约束与肢体间时序协同约束，设计面向四肢联动的奖励项，构建手脚协同的运动生成方案，把该框架可支持的动作类型从单一的下肢 locomotion 扩展到行走取物、跨步搬运等全身协同动作，并检验加入上肢后肌肉激活模式是否仍保持生理合理性。这一改进补全了原框架在动作空间上的功能缺口，而非仅增加动作数量。
	\item 控制算法精度优化。针对原生策略在关节轨迹抖动与稳态跟踪误差上的不足，本课题从观测空间、奖励结构与采样策略三个层面协同改进：以特征选择与归一化降低观测冗余，以平滑惩罚与权重自适应调节缓解多目标奖励之间的冲突，以难样本优先的采样逻辑提升有限预算下的收敛质量。改进目标是在不显著增加网络规模与训练时长的前提下，降低关节角度跟踪误差、抑制轨迹高频抖动，并提升肌肉激活输出的稳定性。
	\item 工程实践创新。完成 KINESIS 的本地化工程改造，解决依赖耦合、日志分散与调试不便等复现难点；将手脚协同模块、精度优化模块与多指标评测脚本封装为可独立调用的组件，形成一套结构清晰、便于二次拓展的工程版本，并提供可复现的训练与评测配置说明，降低后续研究者在该框架上开展实验的门槛。
\end{enumerate}

\section{论文组织架构}
本论文共分为六章，各章内容安排如下。
第 1 章为绪论。阐述课题的研究背景，说明肌肉骨骼运动模仿与 KINESIS 框架的研究意义；梳理物理仿真运动模仿、肌骨建模与仿真、肌肉驱动策略学习以及国内相关研究的发展脉络，评述现有工作的不足；在此基础上明确研究内容、创新点与全文结构。

第 2 章为相关理论与技术基础。介绍人体肌肉骨骼系统的建模原理与仿真求解方式，说明强化学习运动模仿的基本形式与奖励设计，阐述负样本挖掘与混合专家机制的作用机理，并解析KINESIS 的整体架构、模块划分与核心流程，为后续改进提供理论依据。

第 3 章为 KINESIS 框架复现与问题分析。完成实验环境搭建与原始代码部署，复现官方基线实验；通过复现结果定位框架在手脚协同缺失、跟踪精度不足与工程易用性三方面的具体表现，明确改进目标与约束条件。

第 4 章为手脚协同控制与精度优化方法设计。详细给出上肢驱动接入方式、协同约束项的形式与权重设置、观测空间改造方案、奖励函数重构方式以及负样本采样策略的调整，并说明改进后的整体架构、训练流程与实现细节。

第 5 章为实验设计与结果分析。说明数据集、评价指标与实验环境，设置与原生基线的对照实验及关键模块的消融实验，从跟踪精度、协同一致性、平滑度与生理合理性等角度分析结果，并讨论方法的适用范围与局限。

第 6 章为总结与展望。总结全文工作与主要结论，归纳尚未解决的问题，并对肌肉骨骼运动模仿在全身协同、感知噪声建模与实机迁移等方向的拓展作出展望。
